\documentclass[12pt]{article}
\usepackage{amsmath}
\usepackage{amssymb}
\usepackage{natbib}
\usepackage{graphicx}
\usepackage[margin=1in]{geometry}
\usepackage{setspace}
\usepackage{booktabs}
\onehalfspacing
%\doublespacing

\title{United States Then, Europe Now:\\
\large Competing Assumption Plans and the Fiscal Foundations of Federalism}
\author{Thomas J. Sargent}
%\small New York University and Hoover Institution}
\date{}

\begin{document}

\maketitle

\begin{abstract}
This paper reformulates the fiscal history of American federalism by
emphasizing a dimension neglected in \citet{Sargent2012}: the symmetric
competition between state and federal governments in the 1780s to assume
debts and thereby attach creditor interests to their respective levels of
government. Drawing on \citet{Ferguson1961} and \citet{Rothbard2019}, I
characterize the 1780s as a contest between two self-reinforcing
equilibria---a state-dominant equilibrium under the Articles of
Confederation and a federal-dominant equilibrium under the
Constitution---and argue that Hamilton's 1790 assumption plan was the
decisive move in this contest. I then analyze Congress's refusal to
assume state debts in the 1840s as the event that clarified the 1790
settlement was a one-time constitutional bargain, not a permanent
insurance scheme. The reformulation yields sharper lessons for
contemporary debates about European fiscal integration.
\end{abstract}


\section{Introduction}

The fiscal crisis of the 1780s generated competing institutional
responses from state and federal actors, each designed to attach creditor
interests to their level of government. Understanding Hamilton's
ultimate success requires recognizing that he defeated not just inertia
or the inadequacies of the Articles of Confederation, but active
state-level strategies pursuing an alternative fiscal equilibrium. The
battle over assumption was fundamentally about which level of government
would control fiscal capacity and thereby shape the future structure of
American federalism.

Several states pursued aggressive policies to assume Continental
securities held by their residents, creating creditor constituencies
that would naturally oppose transferring taxing authority to a federal
government. These state assumption plans represented a coherent
Anti-Federalist strategy: if successful, they would have preserved the
Articles of Confederation's decentralized allocation of fiscal powers.
Hamilton's federal assumption plan was the Nationalist counter-strategy,
designed to reverse state assumption and consolidate creditor support
behind the federal government.

The paper proceeds as follows. Section~\ref{sec:framework} sets out a
budget-constraint framework that defines what ``competing equilibria''
means in this context. Section~\ref{sec:battle} documents the
state-versus-federal competition for creditor allegiance in the 1780s,
including the politically explosive discrimination debate and the
sinking-fund commitment mechanism. Section~\ref{sec:1840s} analyzes the
1840s state debt crisis and Congress's refusal to extend the 1790
precedent. Section~\ref{sec:europe} draws lessons for contemporary
fiscal unions, particularly the eurozone.


\section{Framework: Government Budget Constraints and Equilibrium
Selection}
\label{sec:framework}

Consider a government that issues one-period debt and whose net-of-interest
surplus follows a stochastic process $\{s_t\}$. A government budget
constraint at date $t$ equates the market value of outstanding debt to
the expected present value of future surpluses:
\begin{equation}
\label{eq:pvbc}
b_t = \sum_{i=0}^{\infty} R^{-(i+1)}\, E_t\!\left(s_{t+1+i}\right),
\end{equation}
where $b_t$ is the real market value of government debt and $R > 1$ is
the gross risk-free interest rate.\footnote{Equation~\eqref{eq:pvbc} is
a present-value budget constraint, not a behavioral equation. It must
hold in any equilibrium in which the government does not default. I use
it here as a diagnostic: a proposed fiscal arrangement is sustainable if
and only if there exists a surplus process $\{s_t\}$ for which
\eqref{eq:pvbc} is satisfied at prices acceptable to creditors.}

The key observation is that the surplus process $\{s_t\}$ is not
exogenous. It depends on fiscal institutions---which level of government
controls which revenue sources, what constraints govern borrowing, and
how creditors' expectations about repayment shape the political
sustainability of taxation. Different institutional arrangements produce
different surplus processes and therefore sustain different levels and
compositions of debt. This is what makes multiple equilibria possible.

Two points about \eqref{eq:pvbc} deserve emphasis. First, the equation
is an \emph{accounting identity} that holds in any non-default
equilibrium; it does not by itself constitute a theory of equilibrium
selection. Second, the interesting economics lies entirely in the
determination of $\{s_t\}$: which taxes are levied, by whom, and with
what political support. I use the framework to organize the historical
narrative, not as a formal model of the coordination game between
creditors and competing levels of government.\footnote{A full
formalization would require a coordination-game setup in which creditors'
expectations about which level of government will service debt are
self-fulfilling: lending to the federal government makes federal revenues
more politically sustainable, which validates the lending. Such a model
would feature strategic complementarities of the type analyzed by
\citet{CooperJohn1988}. I leave this formalization for future work.}


\section{The Battle for Creditor Allegiance: State vs.\ Federal
Assumption}
\label{sec:battle}

\subsection{The Logic of State Assumption}

During the 1780s, several states pursued policies designed to strengthen
state government relative to the Confederation Congress by assuming
Continental securities held by their residents. \citet[pp.~237--239]{Ferguson1961}
documents that Pennsylvania led these efforts, having assumed
approximately \$2~million in Continental debt by 1785. New York and
Massachusetts followed similar strategies. These were not ad hoc
policies but represented a coherent vision of American federalism.

The economic logic mirrored what Hamilton would later employ at the
federal level: creditors who received payment from state governments
would become advocates for preserving state taxing authority. By
converting Continental creditors into state creditors, these states were
implementing a strategy to maintain fiscal sovereignty. If states
successfully assumed most Continental debt owed to their citizens before
any constitutional revision, creditors would have opposed transferring
tariff authority to the federal government---their bonds were already
being serviced by state revenues, and federal assumption would have been
unnecessary from their perspective.

\citet[ch.~48]{Rothbard2019} emphasizes that this represented the
Anti-Federalist economic program: a confederation of sovereign states,
each with its own fiscal capacity, voluntarily cooperating but retaining
ultimate authority over taxation and expenditure. The assumption of
Continental debt by individual states was intended to demonstrate that
state governments could handle the fiscal responsibilities that
Nationalists claimed required a strong central government.

\subsection{The Scale and Mechanics of State Assumption}

\citet{Ferguson1961} provides quantitative evidence of the scope of
state assumption efforts:

\begin{itemize}
\item \textbf{Pennsylvania}: By 1785, Pennsylvania had assumed
  approximately \$2~million in Continental loan certificates and final
  settlement certificates held by Pennsylvania residents. The state
  funded these obligations by issuing state bonds backed by revenues from
  excise taxes and land sales.

\item \textbf{New York}: New York assumed Continental securities and
  funded them through state tariff revenues---the very revenue source
  that would later be transferred to the federal government under the
  Constitution.

\item \textbf{Massachusetts}: Massachusetts attempted extensive
  assumption but struggled with the tax burden this created,
  contributing to Shays' Rebellion in 1786--1787 (see
  Section~\ref{sec:shays} below).

\item \textbf{Virginia and Maryland}: These states paid substantial sums
  on Continental debt through their contributions to Congress, though
  they employed different accounting mechanisms
  \citep[pp.~248--251]{Ferguson1961}.
\end{itemize}

The variation in state approaches is itself revealing. States with
stronger fiscal capacity and more active commerce (Pennsylvania, New
York, Massachusetts) could more credibly pursue assumption strategies.
Southern states faced greater difficulties, which partially explains the
regional divisions at the Constitutional Convention.

\subsection{Two Competing Equilibria}

Two stable institutional outcomes were possible, each corresponding to a
different fiscal structure and a different stochastic process for
government surpluses.

\subsubsection{State-Dominant Equilibrium}

In this equilibrium:
\begin{enumerate}
\item States progressively assume Continental debt owed to their
  residents through the late 1780s.
\item State governments retain tariff authority and other revenue sources
  as specified under the Articles of Confederation.
\item Each state $j$ generates its own surplus process $\{s_{jt}\}$ for
  servicing debt.
\item The federal government remains weak, with insufficient surplus to
  service significant debt.
\item Creditors rationally support state fiscal authority because their
  bonds are serviced from state revenues.
\end{enumerate}

The budget constraint \eqref{eq:pvbc} then applies at the state level:
\begin{equation}
\label{eq:state}
b_{jt} = \sum_{i=0}^{\infty} R^{-(i+1)}\, E_t(s_{j,t+1+i}),
\end{equation}
where $b_{jt}$ is debt issued by state $j$ (including assumed
Continental debt) and $s_{j,t+i}$ is state $j$'s net-of-interest
surplus.

\subsubsection{Federal-Dominant Equilibrium}

In this equilibrium:
\begin{enumerate}
\item The federal government assumes state war debts through Hamilton's
  plan.
\item The federal government acquires tariff authority through the
  Constitution.
\item Federal tariff revenues generate the surplus process
  $\{s_t^{\text{fed}}\}$.
\item State creditors become federal creditors and rationally support
  federal taxation.
\end{enumerate}

Here the budget constraint applies at the federal level:
\begin{equation}
\label{eq:federal}
b_t^{\text{fed}} = \sum_{i=0}^{\infty} R^{-(i+1)}\,
E_t(s_{t+1+i}^{\text{fed}}),
\end{equation}
where $b_t^{\text{fed}}$ includes both original Continental debt and
assumed state debts.

\subsection{The Discrimination Debate and the Politics of Creditor
Allegiance}
\label{sec:discrimination}

Any account of the 1790 assumption centered on creditor allegiance must
confront the politically explosive question of \emph{which} creditors
would benefit. By 1790, a large fraction of Continental and state
securities had been transferred from their original holders---soldiers,
farmers, suppliers who had accepted certificates during the war---to
secondary purchasers, many of them wealthy speculators who had bought at
steep discounts, often 15--25 cents on the dollar
\citep[pp.~271--275]{Ferguson1961}.

Madison proposed \emph{discrimination}: paying current holders the
prevailing market price and distributing the remainder to original
holders. Hamilton opposed discrimination on several grounds.
Operationally, tracing original holders would be prohibitively costly.
But the deeper argument was about capital-market credibility. If the
government discriminated against current holders, it would undermine the
transferability and hence the market value of government
securities---destroying precisely the liquid public debt market that
Hamilton sought to create as the foundation of federal fiscal capacity.

Hamilton prevailed: the Funding Act of 1790 paid current holders at par.
This decision had immense distributional consequences, transferring
wealth to speculators at the expense of original holders. But it also
established that US federal securities were fully transferable claims,
serviced without regard to the identity of the holder. This was
essential for creating a deep and liquid market in federal debt---the
financial foundation of the federal-dominant equilibrium.

The discrimination debate reveals why the competition for creditor
allegiance could not be separated from questions about the character of
public debt itself. The state-dominant equilibrium, with thirteen
separate and relatively illiquid state debt markets, would have looked
very different from Hamilton's vision of a single, nationally traded
federal security.

\subsection{The Sinking Fund as Commitment Device}
\label{sec:sinking}

Hamilton complemented assumption with a sinking fund, established under
the Act of August~12, 1790. The fund was charged with using surplus
revenues to retire debt through open-market purchases, and its
operations were governed by a board of commissioners that included the
Vice President, the Chief Justice, and the Secretary of the Treasury.

The sinking fund served as a commitment device. By earmarking revenues
for debt retirement and placing the fund under a high-profile board,
Hamilton made it politically costly for future Congresses to divert debt
service revenues to other purposes. This was a direct response to the
credibility problem that the Confederation Congress had faced: even when
Congress resolved to service the debt, it lacked institutional mechanisms
to bind itself \citep[pp.~220--225]{Ferguson1961}.

In the language of the framework in Section~\ref{sec:framework}, the
sinking fund was a device to make the surplus process
$\{s_t^{\text{fed}}\}$ more credible---not by increasing surpluses
mechanically, but by constraining the political process that determined
them.

\subsection{The Constitutional Convention as Coordination Device}

The Constitutional Convention and Hamilton's 1790 assumption plan
represented a coordinated move to the federal-dominant equilibrium. But
this required overcoming the partial progress states had already made
toward the alternative equilibrium.

\citet[pp.~289--291]{Ferguson1961} documents the intense political
struggle this created. States that had assumed Continental debt or taxed
their citizens heavily to pay their own debts felt entitled to credit for
those payments. This reflected the reality that these states had already
moved partway toward the state-dominant equilibrium by creating creditor
constituencies attached to state government.

Madison's initial opposition to Hamilton's assumption plan is
illuminating. Virginia had substantially reduced its debt through
taxation, and Madison argued that federal assumption would reward states
that had been fiscally irresponsible while penalizing Virginia's
taxpayers who had already shouldered their burden.

The famous compromise that moved the national capital to the Potomac in
exchange for assumption votes was necessary precisely because Hamilton
was attempting to reverse institutional momentum toward state-level
fiscal consolidation. The accounting system that Hamilton established to
credit states for their prior payments \citep[pp.~293--297]{Ferguson1961}
was essential for overcoming resistance from states that had already
moved toward the state-dominant equilibrium.

\subsection{Why the Federal Equilibrium Prevailed}
\label{sec:shays}

Several factors explain Hamilton's success in establishing the
federal-dominant equilibrium:

\begin{enumerate}
\item \textbf{Shays' Rebellion and distributional conflict}:
  Massachusetts's experience demonstrated not that individual states
  lacked fiscal \emph{capacity}---Massachusetts could and did raise
  revenue---but that the \emph{political economy} of state-level debt
  service was fragile. The burden of hard-money taxes fell
  disproportionately on illiquid western farmers, provoking armed revolt
  in 1786--1787. The rebellion frightened creditors and elites generally:
  even where states could tax, the distributional conflict could make
  debt service politically unsustainable
  \citep[pp.~243--245]{Ferguson1961}. This distinction matters for the
  European analogy, where peripheral eurozone states often have fiscal
  capacity but face intense political resistance to austerity.

\item \textbf{Trade policy coordination problems}: Britain could play
  individual states against each other in trade negotiations. The
  inability to mount a coordinated trade policy revealed a fundamental
  weakness of decentralized fiscal authority.

\item \textbf{Economies of scale in debt service}: Federal consolidation
  allowed a single debt instrument trading in a unified market, reducing
  transaction costs and improving liquidity compared to thirteen separate
  state debt markets.

\item \textbf{Foreign creditor preferences}: Foreign creditors (French
  and Dutch) held primarily Continental rather than state debt and
  preferred dealing with a single federal counterparty
  \citep[pp.~251--253]{Ferguson1961}.

\item \textbf{Speed and coordination}: Hamilton moved quickly after
  ratification. The Acts of August 4 and 5, 1790, were passed less than
  18 months after Washington's inauguration, before state assumption
  plans could advance further.
\end{enumerate}

\citet[ch.~73]{Rothbard2019} offers a more critical interpretation: the
federal equilibrium prevailed because a coalition of creditors, merchants
dependent on trade, and land speculators successfully captured the
constitutional process. Whether one accepts this normative assessment or
not, the account usefully emphasizes that multiple equilibria were
genuinely possible and that the outcome reflected political power rather
than economic inevitability.

\subsection{Implications for Understanding the 1790 Assumption}

Hamilton's 1790 assumption plan was not simply a generous bailout, a
reward for wartime contributions, or a device to create federal fiscal
capacity where none existed. Rather, it was:

\begin{itemize}
\item A strategic move to win the competition for creditor allegiance.
\item A reversal of state-level policies creating alternative creditor
  constituencies.
\item Part of a coordinated institutional transformation
  (Constitution~+ assumption~+ sinking fund) designed to shift to the
  federal-dominant equilibrium.
\item Compensation to states that had moved toward the state-dominant
  equilibrium, inducing them to accept the new federal structure.
\end{itemize}


\section{The 1840s State Debt Crisis}
\label{sec:1840s}

\subsection{The Federal Assumption Debate of the 1840s}

The state debt crisis of the early 1840s must be understood against the
background of the 1790 precedent. Many states had borrowed heavily in the
1830s to finance infrastructure projects---canals, roads, railroads.
During the recession at the end of the 1830s, nine states defaulted on
their debts, and several repudiated them entirely.

European creditors who had purchased these state bonds invoked the 1790
precedent and demanded federal assumption. But they fundamentally
misunderstood what the 1790 assumption had represented. The 1790
assumption settled a constitutional dispute over fiscal federalism by
converting creditors of both Continental and state governments into
federal creditors, as part of a bargain that also transferred fiscal
authority to the federal level. A federal assumption in the 1840s would
have created a permanent guarantee of state debt undertaken for state
purposes, with no constitutional transformation to justify it.

\subsection{The Congressional Debate}

\citet[pp.~165--198]{Ratchford1941} documents the 1840s congressional
debates in detail. Advocates of assumption explicitly cited Hamilton's
1790 plan as precedent. Representative Aaron Vanderpoel of New York
argued that the federal government had a moral obligation to extend the
same treatment to states in the 1840s that it had provided in 1790.

Opponents successfully argued that the situations were fundamentally
different:

\begin{enumerate}
\item Most Continental and state debts assumed in 1790 had been incurred 
  to finance a common national enterprise---the War for Independence.

\item The 1790 assumption was part of a grand constitutional bargain in
  which states surrendered their most important revenue source (tariffs)
  to the federal government.

\item The 1790 assumption was a one-time constitutional settlement, not
  the beginning of a permanent insurance scheme.
\end{enumerate}

By contrast, the debts of the 1830s had been incurred for state-specific
infrastructure projects. There was no national enterprise justifying
federal responsibility. Moreover, many states had borrowed imprudently,
sometimes based on inflated projections of revenues from the very
projects being financed. Federal assumption would create a moral hazard
problem: once states understood that the federal government would
ultimately assume their debts, they would have incentives to borrow
excessively. This is the logic that \citet{KarekenWallace1978} later
formalized in their analysis of deposit insurance, showing that
underpriced government insurance creates incentives for excessive
risk-taking.\footnote{\citeauthor{KarekenWallace1978}'s precise
formulation concerns the incentives of insured banks, not states. But
the structural parallel is close: in both cases, a higher-level
government guarantee of lower-level obligations undermines the
lower-level entity's incentive for fiscal discipline.}

The key congressional opponent was Representative James Garland of
Virginia, who argued that assumption would establish a dangerous
precedent encouraging future state profligacy
\citep[ch.~9]{McGrane1935}.

\subsection{The Decision and Its Consequences}

Congress rejected federal assumption.

\paragraph{Immediate effects.}
The federal government's reputation in European credit markets suffered
temporarily: European creditors did not immediately distinguish between
federal and state creditworthiness, so the state defaults increased
federal borrowing costs. Nine states defaulted (Arkansas, Florida,
Illinois, Indiana, Louisiana, Maryland, Michigan, Mississippi, and
Pennsylvania), and several eventually repudiated portions of their debt
(Arkansas, Florida, Michigan, and Mississippi). Over subsequent years,
creditors learned to distinguish between federal bonds (which had never
defaulted) and state bonds (which carried real default risk).

\paragraph{Long-run institutional effects.}
More importantly, Congress's refusal triggered constitutional reforms at
the state level that persist to this day. \citet{Adams1887} and
\citet{Wallis2005} document how, in response to the defaults and the
absence of federal backstops, voters in more than half the states rewrote
their constitutions in the 1840s and 1850s to impose balanced budget
requirements and debt limitations:

\begin{itemize}
\item Limits on borrowing for operating expenses (as distinct from
  capital projects).
\item Requirements for voter approval of new debt issuance.
\item Debt ceilings as fractions of state property values or tax
  revenues.
\item Balanced operating budgets (though definitions varied across
  states).
\end{itemize}

These constraints remain central to American federalism. Unlike the
federal government, most US states still cannot run sustained operating
deficits.

\subsection{Reinterpreting the 1840s Episode}

The revised narrative helps clarify several aspects of the 1840s crisis:

\begin{enumerate}
\item \textbf{The 1790 assumption was not a general precedent.}
  Congress's 1840s refusal clarified that the 1790 assumption had been a
  one-time settlement of a constitutional dispute, not the establishment
  of permanent federal insurance for state borrowing.

\item \textbf{The refusal prevented moral hazard.} By credibly
  committing not to bail out states that borrowed for state-specific
  purposes, Congress avoided the adverse incentives identified by
  \citet{KarekenWallace1978}.

\item \textbf{The institutional consequences were beneficial.} The state
  constitutional reforms created sustainable constraints on state fiscal
  policy---costly ex post for existing creditors, but establishing clear
  rules for the future.

\item \textbf{Federal fiscal authority was reinforced.} By refusing to
  assume state debts, the federal government paradoxically strengthened
  the clarity of the 1790 settlement: it would service its own debts
  faithfully, but it would not be responsible for state debts.
\end{enumerate}

Hamilton's success in establishing the federal-dominant equilibrium meant
that by the 1830s, there was no ambiguity about institutional
responsibilities. States were sovereign in their own spheres but had no
claim to federal fiscal backing. This reinforces a key lesson:
\emph{federal assumption of lower-level government debts makes sense when
it is part of a constitutional bargain that also transfers fiscal
authority to the federal level; it does not make sense as an ongoing
insurance arrangement where lower-level governments retain fiscal
sovereignty.}


\section{Lessons for Europe and Contemporary Fiscal Unions}
\label{sec:europe}

\subsection{Five Lessons from the American Experience}

The reformulated historical narrative yields five lessons for
contemporary debates about fiscal unions, particularly in Europe.

\paragraph{Lesson 1: Multiple equilibria are possible.}
The United States did not inevitably become a strong fiscal union. A
state-dominant equilibrium under the Articles of Confederation and a
federal-dominant equilibrium under the Constitution were both internally
coherent. The eurozone faces a similar choice: a decentralized
equilibrium with no mutual debt guarantees is feasible, as is a
centralized equilibrium with EU-level fiscal authority and comprehensive
debt mutualization. Intermediate arrangements---a common monetary policy
without fiscal union---may not be sustainable
\citep{DeGrauwe2011}.

\paragraph{Lesson 2: Assumption and fiscal authority must move together.}
The 1790 assumption occurred simultaneously with a transfer of fiscal
authority: Congress acquired exclusive tariff authority while assuming
state debts. The 1840s refusal confirmed the principle from the opposite
direction. European debt mutualization would require a corresponding
transfer of fiscal authority to EU institutions---EU-level taxation,
enforceable fiscal constraints, EU decision-making about public goods,
and reduced member-state sovereignty. Mutualization without these
transfers would replicate the 1840s assumption that Congress rejected,
not the 1790 settlement that established American federalism
\citep[ch.~4]{BrunnermeierJamesLandau2016}.

\paragraph{Lesson 3: Timing and path dependence matter.}
Hamilton moved quickly---the Acts of August 4 and 5, 1790, came only 18
months after Washington's inauguration---before state assumption plans
could create more path dependence toward the state-dominant equilibrium.
Europe's emergency lending facilities (ESM, OMT) have created de facto
fiscal connections among member states without formal fiscal union. The
longer this intermediate state persists, the harder it becomes to move
cleanly in either direction.

\paragraph{Lesson 4: The purpose of the original debts matters.}
The 1790 assumption was justified partly because the debts had financed a
common national enterprise. The 1840s debts financed state-specific
infrastructure and had no comparable collective justification. For
Europe, debts incurred for common purposes (pandemic response, common
defense) may warrant collective responsibility; debts incurred through
national fiscal imprudence have weaker claims to mutualization.

\paragraph{Lesson 5: Credible commitment requires costly demonstrations.}
Congress's 1840s refusal was costly in the short run---European creditors
downgraded all American government bonds. But this short-run cost
purchased long-run credibility and induced states to adopt constitutional
fiscal constraints. Establishing credible no-bailout commitments in
Europe would similarly require allowing defaults when member states
borrow imprudently.


\subsection{Europe's Choices}

The American experience suggests Europe faces a genuine choice between
two coherent institutional arrangements, and that the current
intermediate state is unlikely to be sustainable indefinitely.

\paragraph{Option A: Genuine fiscal union (the 1790 path).}
Comprehensive mutualization of existing member-state debts, transfer of
significant taxation authority to the EU level, EU-level fiscal capacity,
strong constraints on member-state fiscal policies, and a single European
treasury issuing Eurobonds. This is internally coherent but requires
member states to surrender substantial fiscal sovereignty. It implies a
much more integrated political union than currently exists.

\paragraph{Option B: Decentralized confederal system (the reformed
Articles path).}
No mutualization, credible no-bailout commitments demonstrated by
allowing defaults when necessary, member-state constitutional constraints
on borrowing (analogous to post-1840s US state balanced budget
requirements), and regulatory measures to break bank-sovereign doom
loops. This preserves member-state sovereignty but requires accepting
occasional sovereign defaults.

\paragraph{The unstable intermediate.}
Current eurozone arrangements approximate neither Option~A nor Option~B.
The eurozone has a common monetary policy without fiscal union, de facto
bailout mechanisms without explicit mutualization, fiscal rules that are
not fully credible, and no EU-level taxation. The United States spent
less than a decade in a comparable intermediate state (Articles of
Confederation, 1781--1788) before political crisis forced a choice.
Europe has maintained its intermediate state much longer, perhaps because
costs have been distributed gradually---slow growth, high unemployment,
periodic crises---rather than concentrated in a single dramatic event
\citep{BrunnermeierJamesLandau2016, HenningKessler2012}.

\subsection{Important Differences Between the 1780s and Today}

Several differences between the 1780s United States and contemporary
Europe affect the applicability of these lessons. Table~\ref{tab:diffs}
summarizes them.

\begin{table}[htbp]
\centering
\small
\caption{Key differences between the 1780s US and contemporary Europe}
\label{tab:diffs}
\begin{tabular}{@{}p{3.2cm}p{5.5cm}p{5.5cm}@{}}
\toprule
\textbf{Dimension} & \textbf{US, 1780s--1790} &
\textbf{Eurozone, post-1992} \\
\midrule
Initial debt levels & $\sim$40\% of GDP; incurred for a common national
  purpose & Higher \% of GDP; incurred for diverse national purposes \\
Political integration & Shared language, legal traditions, political
  culture, recent revolutionary experience & Greater cultural,
  linguistic, and institutional diversity \\
Labor mobility & High interstate mobility & Low cross-border mobility \\
Government size & Federal expenditures $\sim$1--2\% of GDP & National
  governments $\sim$40--50\% of GDP \\
Social insurance & Minimal; not a source of fiscal heterogeneity &
  Extensive and varying pension, health, and welfare systems \\
Democratic legitimacy & Deliberate constitutional ratification process &
  No comparable constitutional moment for fiscal integration \\
\bottomrule
\end{tabular}
\end{table}

These differences suggest that European fiscal integration would be more
difficult and more consequential than American fiscal integration in
1790. The diversity of European member states and the size of their
governments mean that fiscal union would require more substantial
transfers and harder political compromises.

\subsection{Implications of the Historical Analysis}

The historical analysis implies the following:

\textbf{If the goal is to preserve the eurozone}, European leaders should
move deliberately toward Option~A, following the logic of the 1790
precedent: mutualize existing debts as part of a constitutional bargain,
simultaneously transfer taxation authority to the EU level, and accept
the deep political integration this requires.

\textbf{If deep political integration is unacceptable}, the alternative
is Option~B: strengthen no-bailout commitments credibly, require member
states to adopt constitutional fiscal constraints, break bank-sovereign
doom loops through regulatory means, and accept greater economic
volatility.

\textbf{What the historical analysis counsels against} is continuing
indefinitely in the current intermediate state. The American experience
from the 1780s through the 1840s suggests that maintaining common
monetary institutions without clear resolution of fiscal authority
creates persistent problems: uncertainty about fiscal responsibilities
impairs market functioning, moral hazard accumulates, economic efficiency
suffers, and political legitimacy erodes as citizens cannot identify who
is responsible for fiscal outcomes.


\section{Conclusion}

Incorporating the symmetric competition between state and federal
assumption plans in the 1780s changes how we understand American fiscal
union. Hamilton's 1790 assumption was not a simple bailout or an act of
fiscal generosity. It was the winning move in a strategic competition to
establish the federal-dominant equilibrium over the state-dominant
equilibrium that many Americans preferred. The discrimination debate
and the sinking fund reveal the institutional details through which
Hamilton linked assumption to the creation of a deep and credible federal
debt market. Congress's refusal to assume state debts in the 1840s then
clarified the limits of the 1790 settlement, triggering state-level
institutional reforms that completed the fiscal architecture of American
federalism.

For Europe, the reformulated narrative implies a genuine choice: move
toward deep fiscal integration (the 1790 path) or move toward a more
decentralized system with credible no-bailout commitments (the reformed
Articles path). The current intermediate arrangement is likely not
sustainable indefinitely. The question is whether European leaders will
choose their path deliberately or allow events to choose for them.


\bibliographystyle{plainnat}

\begin{thebibliography}{99}

\bibitem[Adams(1887)]{Adams1887}
Adams, Henry~C. 1887.
\newblock \emph{Public Debts: An Essay in the Science of Finance}.
\newblock New York: Appleton.

\bibitem[Brunnermeier et~al.(2016)Brunnermeier, James, and
  Landau]{BrunnermeierJamesLandau2016}
Brunnermeier, Markus~K., Harold James, and Jean-Pierre Landau. 2016.
\newblock \emph{The Euro and the Battle of Ideas}.
\newblock Princeton: Princeton University Press.

\bibitem[Cooper and John(1988)]{CooperJohn1988}
Cooper, Russell, and Andrew John. 1988.
\newblock ``Coordinating Coordination Failures in Keynesian Models.''
\newblock \emph{Quarterly Journal of Economics} 103 (3): 441--63.

\bibitem[De~Grauwe(2011)]{DeGrauwe2011}
De~Grauwe, Paul. 2011.
\newblock ``The Governance of a Fragile Eurozone.''
\newblock \emph{CEPS Working Document} No.~346.

\bibitem[Ferguson(1961)]{Ferguson1961}
Ferguson, E.~James. 1961.
\newblock \emph{The Power of the Purse: A History of American Public
  Finance, 1776--1790}.
\newblock Chapel Hill: University of North Carolina Press.

\bibitem[Henning and Kessler(2012)]{HenningKessler2012}
Henning, C.~Randall, and Martin Kessler. 2012.
\newblock ``Fiscal Federalism: US History for Architects of Europe's
  Fiscal Union.''
\newblock \emph{Bruegel Essay and Lecture Series}.

\bibitem[Kareken and Wallace(1978)]{KarekenWallace1978}
Kareken, John~H., and Neil Wallace. 1978.
\newblock ``Deposit Insurance and Bank Regulation: A Partial-Equilibrium
  Exposition.''
\newblock \emph{Journal of Business} 51 (3): 413--38.

\bibitem[McGrane(1935)]{McGrane1935}
McGrane, Reginald~C. 1935.
\newblock \emph{Foreign Bondholders and American State Debts}.
\newblock New York: Macmillan.

\bibitem[Ratchford(1941)]{Ratchford1941}
Ratchford, B.~U. 1941.
\newblock \emph{American State Debts}.
\newblock Durham, NC: Duke University Press.

\bibitem[Rothbard(2019)]{Rothbard2019}
Rothbard, Murray~N. 2019.
\newblock \emph{Conceived in Liberty, Volume 5: The New Republic,
  1784--1791}.
\newblock Auburn, AL: Mises Institute.

\bibitem[Sargent(2012)]{Sargent2012}
Sargent, Thomas~J. 2012.
\newblock ``United States Then, Europe Now.''
\newblock \emph{Journal of Political Economy} 120 (1): 1--40.

\bibitem[Wallis(2005)]{Wallis2005}
Wallis, John~Joseph. 2005.
\newblock ``Constitutions, Corporations, and Corruption: American States
  and Constitutional Change, 1842 to 1852.''
\newblock \emph{Journal of Economic History} 65 (1): 211--56.

\end{thebibliography}

\end{document}
